% Modificaciones generales para el documento
% Por E.Vidal
% v0.1 09/24.
%%%%%%%%%%%%%%%%%%%%%%%%%%%%%%%%%%%%%%%%%%%%%%%%%%%%%%%%%%%%%%%%%%%%%%%%
%% Packages related to general page setup
%\usepackage[hmargin={2.0cm,2.0cm},height=24cm]{geometry}
\usepackage{lipsum}         % Eliminar este paquet#e solo es para generar texto de muestra.
\usepackage[onehalfspacing]{setspace}   % distance between lines 1-1/2-lines
\usepackage[utf8]{inputenc} % caracteres especiales como la ñ
\usepackage[spanish]{babel} % caracteres del español
\spanishdecimal{.}
\usepackage[T1]{fontenc}    % correcta representacion de varias tildes/acentos, p.ej. en alemán 'ß'
\usepackage{csquotes}
\usepackage[backend=biber,natbib=true,style=apa]{biblatex}
\setlength\bibitemsep{0.6\baselineskip} % Adds spaces between references entries
\addbibresource{style/bibliografia.bib}
%% Global change of the font
\usepackage{lmodern}           % latin modern typography
\usepackage{titling}           %to ensure image on cover (caratula)
% commands for including the picture
\newcommand{\titlepicture}[2][]{%
  \renewcommand\placetitlepicture{%
    \includegraphics[#1]{#2}\par\medskip
  }%
}
\linespread{1.25}
\newcommand{\placetitlepicture}{} % initialization

\usepackage{scrhack} % get rid of warning "\float@addtolists detected"
\usepackage[automark]{scrlayer-scrpage}  
\pagestyle{scrheadings}		           % find explanation in scrguide.pdf 
\clearpairofpagestyles
% To have: "Chapter X. ..."
\renewcommand*{\chapterformat}{\mbox{\chapapp{\nobreakspace}\thechapter\autodot\IfUsePrefixLine{}{\enskip}}}
\renewcommand*{\chaptermarkformat}{\chapapp~\thechapter\autodot\enskip}
%\automark[subsection]{section}        % section im Header with Nr.  -- still not working.
% oneside optionen Kopfzeile:
\ihead[]{\headmark} 
\chead[]{}
\ohead[]{\namedesautors}
% oneside optionen Fußzeile:
\ifoot[]{}
\cfoot[\thepage]{}
\ofoot[]{\thepage}
%\addtokomafont{pagehead}{\upshape}% header upshape instead italic
%
% twoside optionen Kopfzeile:
%\lehead[Inhalt plain.scrheadings]{Inhalt scrheadings}
%\cehead[Inhalt plain.scrheadings]{Inhalt scrheadings}
%\rehead[Inhalt plain.scrheadings]{Inhalt scrheadings}
%\lohead[Inhalt plain.scrheadings]{Inhalt scrheadings}
%\cohead[Inhalt plain.scrheadings]{Inhalt scrheadings}
%\rohead[Inhalt plain.scrheadings]{Inhalt scrheadings}
% twoside optionen Fußzeile:
%\lefoot[Inhalt plain.scrheadings]{Inhalt scrheadings}
%\cefoot[Inhalt plain.scrheadings]{Inhalt scrheadings}
%\refoot[Inhalt plain.scrheadings]{Inhalt scrheadings}
%\lofoot[Inhalt plain.scrheadings]{Inhalt scrheadings}
%\cofoot[Inhalt plain.scrheadings]{Inhalt scrheadings}
%\rofoot[Inhalt plain.scrheadings]{Inhalt scrheadings}
\addtokomafont{chapter}{\fontsize{14}{12}\selectfont}
\addtokomafont{section}{\large}
\addtokomafont{subsection}{\large}
\addtokomafont{subsubsection}{\large}
%%%%%%%%%%%%%%%%%%%%%%%%%%%%%%%%%%%%%%%%%%%%%%%%%%%%%%%%%%%%%%%%%%%%%%%%%
%% Table of Contents (ToC) re-Declaration
\newcommand\partentrynumberformat[1]{\partname\ #1}
\RedeclareSectionCommand[
  tocentrynumberformat=\partentrynumberformat,
  tocnumwidth=4em
]{part}

\newcommand\chapterentrynumberformat[1]{\chapapp\ #1}
\RedeclareSectionCommand[
    tocentrynumberformat=\chapterentrynumberformat,
    tocnumwidth=5.5em
]{chapter}

\newcommand*{\appendixmore}{
  \renewcommand*{\chapterformat}{\appendixname~{\nobreakspace}\thechapter\enskip}%
  \renewcommand*{\chaptermarkformat}{\appendixname~{\nobreakspace}\thechapter\enskip}%
  \addtocontents{toc}{\protect\appendixtocentry}%
}

\newcommand*{\appendixtocentry}{%
  \DeclareTOCStyleEntry[
    entrynumberformat=\appendixprefixfixintoc,
    dynnumwidth
  ]{default}{chapter}%
}
\newcommand{\appendixprefixfixintoc}[1]{%
  \def\autodot{}%
  \appendixname~#1%
}

\RedeclareSectionCommand[tocindent=5.5em]{section}
\RedeclareSectionCommand[tocindent=7.8em]{subsection}
\RedeclareSectionCommand[tocindent=11em,toclevel=3]{subsubsection} % Asegura que aparezca en el TOC
\setcounter{tocdepth}{3}  % Muestra hasta subsubsection en el índice
\setcounter{secnumdepth}{3}
 % Numera hasta subsubsection

%%%%%%%%%%%%%%%%%%%%%%%%%%%%%%%%%%%%%%%%%%%%%%%%%%%%%%%%%%%% set indentation
\def\EspaciadoInicioParrafo{20 pt}
\def\EspaciadoEntreParrafos{0 pt}
\def\EspacioEnBlancoLastLine{0pt plus 1fil}         %white space required at the end of the last line
\setparsizes{\EspaciadoInicioParrafo}{\EspaciadoEntreParrafos}{\EspacioEnBlancoLastLine}
%%%%%%%%%%%%%%%%%%%%%%%%%%%%%%%%%%%%%%%%%%%%%%%%%%%%%%%%%%%%%%%%%%%%%%%%
%% Packages related to AUTHOR
%%----------------------------------------------------------------------------
%%-- Info about the work and author
%%
%
%\newcommand{\arbeitGeneralDescriptionPre}{Trabajo de Investigaci\'on para obtener el grado de }
%\newcommand{\arbeitGeneralDescriptionPost}{EN INGENIER\'IA MECATR\'ONICA}
\newcommand{\arbeitGeneralDescriptionPre}{Trabajo de Fin de Carrera}
\newcommand{\arbeitGeneralDescriptionPost}{1/2}                           % 1 ó 2 para TFC1 o TFC2 respectivamente
\newcommand{\artderarbeit}{BACHILLER}                                   % Bachiller o Magister
\newcommand{\titelderarbeit}{SISTEMA MECATR\'ONICO PARA XXXX\\[0.2ex] AUTOMATIZADO DE YYYYY\\[1ex] CON EL USO DE ZZZZZ}    % Title of your work
\newcommand{\namedesautors}{Grace Stefani Plantilla Muestra}            % Nombre tesista (nombres apellidos)
\newcommand{\matrikelnummer}{20212123}                                  % Codigo de estudiante
\newcommand{\studiengangdesautors}{Mecatr\'onica}                       % study degree programme here
\newcommand{\namedesfachgebiets}{Mecatr\'onica}                         % Nombre Facultad
\newcommand{\namedesprofessors}{Mag. Ing. Enrique Vidal}                % Nombre Asesor
\newcommand{\namedeswissenschaftlichenbetreuers}{KEINER}                % Nombre 2ndo Asesor, ingresar 'KEINER' si no tiene.
\newcommand{\namedesbetrieblichenbetreuers}{KEINER}                     % No es necesario llenar dejarlo con 'KEINER'
\newcommand{\abgabedatum}{\today}                                       % Datum of delivery.
\newcommand{\nameUniversity}{Pontificia Universidad Cat\'olica del Per\'u}  % Nombre Universidad
%%%%%%%%%%%%%%%%%%%%%%%%%%%%%%%%%%%%%%%%%%%%%%%%%%%%%%%%%%%%%%%%%%%%%%%%
%% Packages related to general page setup
\usepackage[table,xcdraw,dvipsnames]{xcolor}%para definir colores
\AfterFile{hdvips.def}{\hypersetup{linktocpage}}
\AfterFile{hypertex.def}{\hypersetup{linktocpage}}
\usepackage{hyperref}    % enable hyperrefs
\PassOptionsToPackage{unicode}{hyperref}
\PassOptionsToPackage{naturalnames}{hyperref}
\pdfstringdefDisableCommands{%
  \let\HyPsd@CatcodeWarning\@gobble
  \def\\{}%
  \def\texttt#1{<#1>}%
  \let\@ac\relax%
}
%-- hypersetup
\hypersetup{
	pdftitle     = {\titelderarbeit},
	pdfsubject   = {\artderarbeit},
	pdfauthor    = {\namedesautors},
	pdfkeywords  = {\artderarbeit; \namedesfachgebiets; {\nameUniversity};},
	colorlinks   = true,
	linkcolor    = white!0!black,
	urlcolor     = white!0!black,
	citecolor    = black,
	pdfstartview = {Fit}
}
\usepackage{pdfpages}

\usepackage{tikz}
%\usetikzlibrary{arrows.meta} % Para estilos de punta de flecha


\usepackage{pgfplots}
\usepackage{pgfplotstable}
\pgfplotsset{compat=1.18}


\usepackage{xcolor}

\usepackage[printonlyused]{acronym} %for acronyms
\makeatletter
\newif\if@in@acrolist
\AtBeginEnvironment{acronym}{\@in@acrolisttrue}
\newrobustcmd{\LU}[2]{\if@in@acrolist#1\else#2\fi}

\newcommand{\ACF}[1]{{\@in@acrolisttrue\acf{#1}}}
\makeatother
\renewcaptionname{spanish}{\listtablename}{\'Indice de tablas}
\renewcaptionname{spanish}{\listfigurename}{\'Indice de figuras}
\renewcaptionname{spanish}{\tablename}{Tabla} %used on longtblr
\renewcaptionname{spanish}{\appendixname}{Anexo} 
\spanishlcroman
%%%%%%%%%%%%%%%%%%%%%%%%%%%%%%%%%%%%%%%%%%%%%%%%%%%%%%%%%%%%%%%%%%%%%%%%
%% Common Packages
\usepackage{graphicx} %para insertar imagenes
%\graphicspath{{./bilder/}}   % Path to images (optional)
\usepackage{epsfig}           % Package to link postscript Graphics
\pdfminorversion=7
%\usepackage{float}           % Macro which helps to correctly display the position of graphics

\usepackage{mdframed}  % Para crear recuadros personalizados
\usepackage{amsmath} 
\usepackage{amssymb}
\usepackage{amsthm}
\usepackage{siunitx}
\usepackage{xparse}

% "Desactiva" el formateo numérico de \SI y solo pega número + unidad:
\RenewDocumentCommand{\SI}{m m}{#1\,\si{#2}}

% (Opcional) Para rangos:
\RenewDocumentCommand{\SIrange}{m m m}{#1--#2\,\si{#3}}

% (Opcional) Si usas \num o \numrange y también te cambian el separador:
\RenewDocumentCommand{\num}{m}{#1}
\RenewDocumentCommand{\numrange}{m m}{#1--#2}

% \sisetup{
%   output-decimal-marker = {.}
% }
% \sisetup{locale = US}
\usepackage[shortcuts]{extdash}

%insertar tablas
\usepackage{codehigh} % Para usar \fakeverb (similar a \verb!!) dentro de tablas
\usepackage{tabularx} % Asegúrate de tener cargado este paquete / util para centrar imagenes
\usepackage{array}    % Para mejorar el control del ancho de las columnas
\usepackage{multicol}
\usepackage{multirow}
\usepackage{colortbl}
\usepackage{tabularray}
\setlength\bibitemsep{0.6\baselineskip} % Adds spaces between references entries
\usepackage{hhline} % For \hhline in tables
\DefTblrTemplate{contfoot-text}{es}{Continúa en la siguiente página}
\DefTblrTemplate{conthead-text}{es}{(Continúa)}
% Para no declarar el diseño de tabla utilizar esta línea de comando
% \DefTblrTemplate{contfoot-text}{default}{Continúa en la siguiente página}
% \DefTblrTemplate{conthead-text}{default}{(Continúa)}
\SetTblrTemplate{contfoot-text}{default}
\SetTblrTemplate{conthead-text}{default}
\UseTblrLibrary{amsmath}
%See https://tex.stackexchange.com/questions/621766/figures-with-caption-inside-tabularray-tblr-tables
%% also: https://github.com/lvjr/tabularray/discussions/47
\UseTblrLibrary{counter,varwidth} %To fix some counter in images within tables
\NewTblrTheme{pucp}{
\SetTblrStyle{firsthead}{}
\SetTblrTemplate{contfoot-text}{es}
\SetTblrTemplate{conthead-text}{es}
% Para no declarar el diseño de tabla utilizar esta línea de comando
% \SetTblrTemplate{contfoot-text}{normal}
% \SetTblrTemplate{conthead-text}{normal}
}

\usepackage{bigstrut}
\usepackage{booktabs}
\usepackage{float}
\usepackage{caption}
\usepackage{longtable}
\captionsetup[table]{name = Tabla}
\usepackage{subcaption}
%par formato APA
%\captionsetup{justification=raggedright, singlelinecheck=false} 
\usepackage{placeins}  % Colocar en el preámbulo
\KOMAoptions{DIV=last}  %recalculates the dimensions, siempre al final.
%User defined commands
%Definiciones para uso en ecuaciones
\newcommand{\mbbb}[2]{\mathbb{R}^{#1\times{#2}}}
%definiciones de separacion de codigo
\def\MPWxSMALLSxLISTING{0.55}

%\def\EspaciadoInicioParrafo{1.0}
%\hspaceEspaciadoInicioParrafo
%\newcommand{\hEIPar}{\hspace{\EspaciadoInicioParrafo cm}}

%definiones de separacion de tablas 


%\begin{minipage}{0.30\textwidth}
%\vspaceTextToFigure % =\vspace{0.1cm}
%insertar imagen
%\captionof{figura}{Red LED selection}
%\end{minipage}
%\begin{minipage}{0.70\textwidth}
%\vspaceTextToFigure % =\vspace{0.1cm}
%\end{minipage}
%conitnua texto

\usepackage{listings,lstautogobble}
\usepackage{listingsutf8}   % comentarios con acentos en UTF-8

% Estilo base (evita overflow)
\lstdefinestyle{codebase}{
  inputencoding=utf8,
  backgroundcolor=\color{codebg},
  frame=single, rulecolor=\color{codeborder}, framerule=0.6pt, framesep=6pt,
  numbers=left, numberstyle=\tiny, numbersep=8pt,
  xleftmargin=1.5em, xrightmargin=0em,
  basicstyle=\ttfamily\footnotesize,
  keywordstyle=\bfseries\color{codekw},
  commentstyle=\itshape\color{codecm},
  stringstyle=\color{codestr},
  identifierstyle=\color{codeid},
  showstringspaces=false, upquote=true, tabsize=2, keepspaces=true,
  breaklines=true,            % <--- quiebra líneas largas
  breakatwhitespace=false,    % permite cortar dentro de tokens si es necesario
  columns=fullflexible,       % mejora el ajuste para evitar desbordes
  postbreak=\mbox{\textcolor{codecm}{\(\hookrightarrow\)}\space}, % marca de continuación
  captionpos=b, aboveskip=1ex, belowskip=1ex
}

% Estilo específico para C
\lstdefinestyle{cstyle}{
  style=codebase,
  language=C,
  morekeywords={uint8_t,uint16_t,uint32_t,bool,true,false}
}

% MATLAB/Octave (listings suele traer 'Matlab'; si no, usa 'Octave')
\lstdefinestyle{matlabstyle}{
  style=codebase,
  language=Matlab        % o: language=Octave
}

%  Mapea caracteres españoles en comentarios si hiciera falta:
\lstset{
  literate=
   {á}{{\'a}}1 {é}{{\'e}}1 {í}{{\'i}}1 {ó}{{\'o}}1 {ú}{{\'u}}1
   {Á}{{\'A}}1 {É}{{\'E}}1 {Í}{{\'I}}1 {Ó}{{\'O}}1 {Ú}{{\'U}}1
   {ñ}{{\~n}}1 {Ñ}{{\~N}}1 {°}{{$^\circ$}}1
}
\renewcommand{\lstlistoflistings}{\begingroup
  \tocfile{\lstlistlistingname}{lol}
\endgroup}
% Nombres en español
\renewcommand{\lstlistingname}{Código}
\renewcommand{\lstlistlistingname}{Índice de códigos}

%\ExplSyntaxOff
\usepackage{pdflscape}
\usepackage[table,xcdraw]{xcolor}